\chapter{State of the Art}\label{ch:state-of-the-art}


\section{The CS-Aware-Next Project}\label{sec:the-cs-aware-next-project}
The CS-Aware-Next project{\footnote{\href{https://www.cs-aware-next.eu/}{https://www.cs-aware-next.eu/}}} is an innovative European initiative dedicated to advancing cybersecurity by integrating state-of-the-art artificial intelligence and data analytics techniques. This collaborative project brings together experts from academia, industry, and government to develop robust, scalable solutions to detect, analyze, and mitigate cyber threats in real-time.

The primary objectives of CS-Aware-Next are to develop advanced machine learning and deep learning algorithms to automatically identify and classify cybersecurity threats from a variety of heterogeneous data sources, e.g., textual data and network traffic, and consequently, to leverage modern natural language processing techniques, such as state-of-the-art embedding models and encoder architectures, to transform unstructured cyber intelligence into informative, dense vector representations that support multi-modal data fusion. The CS-Aware-Next project also includes creating scalable, real-time monitoring systems that can efficiently process large volumes of data with the property to adapt to evolving threat landscapes. For this reason, this project integrates interdisciplinary insights from computer science, statistics, and cybersecurity to address challenges such as class imbalance, label dependencies, and the high dimensionality of cyber data. The project employs a range of methodologies closely related to those discussed in this thesis. These include using transformer-based models (e.g., BERT and its variants), advanced ensemble methods (such as Random Forests and XGBoost), and state-of-the-art natural language processing techniques to extract and classify security-relevant information.

The techniques presented in this thesis directly apply to the goals of CS-Aware-Next. The methods discussed enable practical analysis and automated detection of cyber threats by translating unstructured text into high-dimensional semantic representations, which is at the heart of the CS-Aware-Next project. This thesis contributes to the broader research landscape by trying to demonstrate how modern deep learning and NLP techniques can be leveraged to extract meaningful features from complex, high-dimensional data to provide insights into the challenges and solutions associated with multi-label classification and representation learning, both critical for processing cybersecurity intelligence, and to highlight the potential for AI-driven approaches to enhance real-time threat detection and improve cybersecurity defenses.

The CS-Aware-Next project represents a significant step towards developing autonomous cybersecurity systems that can anticipate and respond to emerging threats with minimal human intervention. Its focus on integrating advanced AI with cybersecurity practices aims to enhance the detection of sophisticated cyberattacks by enabling the early identification of threat patterns in large-scale data and promoting interdisciplinary collaboration research.


\section{AI-based Analysis}\label{sec:ai-based-analysis}
This chapter presents the essential tools for understanding the methodologies and results of the thesis. Artificial intelligence is introduced initially, along with its significant subfields, i.e., machine learning and deep learning. Then, the natural language processing section is presented as a consequence of using deep learning techniques on text data and, because of it, the popular tools for embeddings and encoders. Presented how many embeddings, and so the encoders, are significant tools in translating text data into numerical vectors for the subsequent algorithms; the discussion proceeds with the presentation of the algorithms used in this project for achieving its aim.

\textbf{Artificial Intelligence (AI)} is a broad field focused on creating intelligent entities or machines that can act effectively and safely in novel situations. It involves both understanding and building intelligent systems. AI is the study of agents that perceive their environment and perform actions. The goal of AI has been defined as creating systems that maximize expected utility, but this is evolving to include systems that learn what to maximize and function appropriately, even when uncertain about the objective. Learning is emphasized to construct competent systems and extend the designer's reach into unknown environments. AI includes various subfields, from general areas like machine learning, deep learning, and natural language processing. \cite{russellnorvig2009}

\textbf{Machine Learning (ML)} allows computers to improve their performance through experience. Instead of explicit programming for every situation, machine learning involves a system that observes data, builds a model based on that data, and then uses the model to make predictions or decisions. The generated model serves as a hypothesis and a problem-solving tool. ML becomes valuable when designers cannot anticipate all potential future situations or do not know how to program a solution directly. In the ML subfield, the branch of Supervised Learning consists of applications where the training data comprises examples of the input vectors along with their corresponding target vectors. The aim is to assign each input vector to one of a finite number of discrete categories called classification problems. \cite{bishop2006}

A \textbf{Neural Network (NN)} is an engineered system inspired by the biological brain and used for machine learning. It can be considered a function approximation machine designed to achieve statistical generalization, occasionally drawing insights from the brain. It is also referred to as an artificial neural network, ANN. In a feedforward network, information flows through the evaluated function, from input $x$ to intermediate computations and finally to the output $y$. Feedforward neural networks can be extended to include feedback connections, at which point they are called recurrent neural networks. Designing the network's architecture involves deciding how many layers it should contain, how these layers should be connected, and how many units should be in each. \cite{goodfellowbengiocourville2016}

\textbf{Deep Learning (DL)} is a type of ML that uses ANN with multiple layers (depth) to analyze data and extract patterns. Deep learning expands on the basic principles of neural networks by increasing the depth and complexity of the models. ``Deep learning'' emphasizes the ability to train deeper neural networks than previously possible, focusing on the theoretical importance of depth. Modern deep learning emphasizes learning multiple levels of composition and can be applied in machine learning frameworks that are not necessarily neurally inspired. Deep learning was developed to address the limitations of traditional machine learning algorithms in generalizing complex AI tasks. The popularity and usefulness of deep learning have been driven by more powerful computers, larger datasets, and techniques to train deeper networks. \cite{goodfellowbengiocourville2016}

A \textbf{Transformer}~\cite{vaswani2023} is a neural network architecture that relies on attention mechanisms to model relationships in data, like text or code. The attention mechanism is central to the Transformer's design. It allows the model to weigh the importance of different input parts when processing information. It helps the model learn how tokens relate to each other, even over long distances. Transformers use self-attention, allowing each word to attend to all other words in the input sentence to understand the context better. Transformers often employ multiple attention heads, allowing the model to capture different types of relationships in the data. Each head may specialize in different linguistic relationships or patterns. \cite{jurafsky2025}

\subsection{Natural Language Processing}\label{subsec:natural-language-processing}
The \textbf{Natural Language Processing (NLP)} literature aims to create autonomous models in the feature extraction step by training them to be self-supervised, avoiding the need for hand-labeled supervision signals. The resulting model will be the embedding model, or embedder, for texts, and it is noticeable that it will compute dense semantic vectors for each sentence in input so that words related in semantics will end close in the embedding space representation, in which each point is associated to a word meaning. \cite{jurafsky2025}

\subsubsection{The Embeddings}\label{subsubsec:the-embeddings}
An \textbf{embedding} represents data in a vector space. It involves mapping words (or other data types) into vectors of real numbers. These vectors capture semantic and syntactic information about the words, allowing mathematical operations to model relationships between them. An embedding represents a word as a point in a multidimensional semantic space. There exist two types of embeddings: static embeddings, which assign a fixed vector to each word in the vocabulary (e.g., word2vec, GloVe), and contextual embeddings, which represent the meaning of a word in its specific context, so the vector for a word changes depending on the sentence it appears in (e.g., BERT). \cite{jurafsky2025}

\textbf{Word2vec}~\cite{mikolovchencorradodean2013} is a method for computing embeddings. It is fast, efficient to train, and easily accessible online with code and pre-trained embeddings. Word2vec embeddings are static embeddings, which means that the method learns one fixed embedding for each word in the vocabulary. Instead of counting how often each word occurs near a target word, word2vec trains a classifier on a binary prediction task: ``Is word $w$ likely to show up near \texttt{target}?'' The learned classifier weights are then taken as the word embeddings. \cite{jurafsky2025}

\textbf{Positive Pointwise Mutual Information (PPMI)} is an alternative weighting function to TF-IDF. It is used for term-term matrices when the vector dimensions correspond to words rather than documents. PPMI is based on the intuition that the best way to weigh the association between two words is to determine how much more the two words co-occur in a corpus than would be expected by chance. \cite{jurafsky2025}

\textbf{GloVe (Global Vectors)}~\cite{penningtonsochermanning2014} is a widely used static embedding model that captures global corpus statistics. It is based on ratios of probabilities from the word-word co-occurrence matrix, combining the intuitions of count-based models like PPMI while also capturing the linear structures used by methods like word2vec. It is a static embedding since it learns one fixed embedding for each word in the vocabulary to capture the global corpus statistics. \cite{jurafsky2025}

\subsubsection{The Encoders}\label{subsubsec:the-encoders}
\textbf{BERT (Bidirectional Encoder Representations from Transformers)}~\cite{devlin2019} is a bidirectional transformer encoder model trained via Masked Language Modeling, MLM. Unlike causal (left-to-right) transformers, bidirectional encoders use self-attention to map input embeddings to output embeddings, contextualized using information from the entire input sequence. BERT produces dynamic or contextualized embeddings where the vector for each word is different in different contexts. During training, a percentage of tokens are randomly chosen to be manipulated. Thus, the model masks the selected words in the middle sentences and guesses them based on their surrounding words. This approach allows the model to consider both left and right context. The output vector from the model's final layer represents the token's meaning in its context. These contextual embeddings can then be used for tasks requiring the meaning of tokens or words. \cite{jurafsky2025}

\textbf{DistilBERT} is a distilled version of BERT. It is a smaller, faster, cheaper, and lighter general-purpose language representation model that can be fine-tuned with good performance on various tasks. DistilBERT has 40{\%} fewer parameters and is 60{\%} faster than BERT while retaining 97{\%} of BERT's language understanding capabilities. It is pre-trained through knowledge distillation, using a triple loss that combines language modeling, distillation, and cosine-distance losses. DistilBERT is trained on the same corpus as the original BERT model, a concatenation of English Wikipedia and Toronto Book Corpus. \cite{sanh2020}

\textbf{RoBERTa (Robustly Optimized BERT Approach)} improves upon the original BERT pre-training procedure. Unlike the original BERT implementation, which performed masking once during data preprocessing, RoBERTa uses dynamic masking. Dynamic masking is either comparable to or slightly better than static masking. RoBERTa removes the next sentence prediction objective used in BERT\@. It utilizes a more extensive byte-level BPE vocabulary containing 50K sub-word units without additional preprocessing or tokenization of the input. \cite{liu2019}

\textbf{Sentence-BERT (SBERT)} is a modification of the BERT network that uses siamese and triplet network structures to derive semantically meaningful sentence embeddings that can be compared using cosine-similarity. SBERT is designed to map sentences to a vector space where semantically similar sentences are close to each other. It reduces the computational effort required to find the most similar pair in a collection of sentences, making it suitable for semantic similarity search and clustering. For example, finding the most similar pair in a collection of 10.000 sentences is reduced from 65 hours with BERT/RoBERTa to about 5 seconds with SBERT\@. It is fine-tuned on Natural Language Inference (NLI) data. SBERT outperforms other state-of-the-art sentence embedding methods on Semantic Textual Similarity (STS) tasks and transfer learning tasks. \cite{reimersgurevych2019}

\textbf{ATT{\&}CK BERT} is an attack vector embedding model. Its primary function is to create a semantically meaningful embedding where attack vectors with similar meanings are close in the embedding space, indicated by high cosine similarity. It is a transformer model fine-tuned using the SIAMESE network over the ATT{\&}CK Matrix data. It is trained to maximize the cosine similarity between pairs of attack vectors extracted from the same ATT{\&}CK technique and to minimize the cosine similarity between attack vectors, each extracted from a different ATT{\&}CK technique. Compared to models like the original SBERT trained on general datasets, ATT{\&}CK BERT is specifically fine-tuned on the ATT{\&}CK matrix. This enables it to extract embeddings that better represent attack vectors in the cybersecurity domain. \cite{basel2023}

\textbf{ATT{\&}CK} is a knowledge base developed by MITRE that categorizes attack techniques that adversaries have been observed using in the real world. ATT{\&}CK categorizes attack techniques into different tactics, where each tactic represents a goal that the attacker tries to achieve, such as initial access, execution, or defense evasion. A tactic contains multiple techniques, including a textual description of the attacker's behaviors and examples of real-world uses of the method by known malicious groups and software. Each technique also contains detection and mitigation recommendations to aid security teams in countering the technique. Thanks to the high-level attack categorization and practical defense recommendation offered by ATT{\&}CK, it is a valuable resource for security analysts to secure their systems and networks. \cite{basel2023}


\section{State-of-the-Art Text Classifiers}\label{sec:state-of-the-art-text-classifiers}
\textbf{Text classification} is the task of assigning a category to a text or document. It involves taking an observation, extracting useful features, and classifying it into discrete classes. Text classification algorithms map an input $x$ to a predicted class $y$ from a fixed set of output classes $Y = \left\{y_1, y_2, \dots, y_M\right\}$. In text classification, $c$ often represents ``class'' and $d$ represents ``document''. Text classification is commonly done via supervised machine learning. The goal is to learn a classifier to map a new document $d$ to its correct class $c \in C$, where $C$ is a set of helpful document classes. A probabilistic classifier also provides the probability of the observation being in a particular class. \cite{jurafsky2025}

\subsection{Binary Classification Algorithms}\label{subsec:binary-classification}
\textbf{Logistic regression} is a classification model where the posterior probability of a class can be written as a logistic sigmoid acting on a linear function of the feature vector. It is a generalized linear model used for classification, not regression, despite its name. The goal is to predict discrete class labels or posterior probabilities in the $\left(0, 1\right)$ range. \cite{bishop2006}

A \textbf{Naïve Bayes classifier} is a probabilistic classifier that simplifies how features interact. It is a type of Bayesian classifier that uses Bayes' theorem but makes a ``naïve'' assumption of conditional independence among the features given the class. Naïve Bayes models assume that the value of one feature is independent of the value of any other feature, conditioned on the class. This assumption simplifies the model and makes it computationally efficient, but it is rarely true in real-world applications. A Naïve Bayes classifier commonly represents a text document as a bag of words in text classification. This means it ignores word order and grammatical structure, considering only word frequencies within the document. Even if its conditional independence assumption is not precisely satisfied, the Naïve Bayes model may still give good classification performance because decision boundaries can be insensitive to some details in the class-conditional densities. This technique is helpful when the dimensionality $D$ of the input space is high, making density estimation in the full $D$-dimensional space more challenging. On the other hand, the conditional independence assumption may lead to poor representations of the class-conditional densities. \cite{bishop2006}

A \textbf{Random Forest} is an ensemble learning method that builds a prediction model by combining the strengths of a collection of simpler base models such as trees. The idea is to improve bagging's variance reduction by reducing the correlation between the trees without increasing the variance too much. It grows multiple trees on bootstrapped data, and each tree gives a classification result. Before each split, the random forest selects $m \leq p$ of the input variables at random as candidates for splitting, where $p$ represents the total number of input variables available in the dataset. Then, the random forest classifies using a majority vote. For classification, a random forest obtains a class vote from each tree and then classifies using a majority vote. \cite{james2013}

\textbf{XGBoost (eXtreme Gradient Boosting)} is a scalable, end-to-end tree boosting system widely used in machine learning. It implements gradient boosting, a machine learning technique that combines the outputs of many ``weak'' classifiers to produce a powerful ``committee''. Gradient boosting involves creating an additive model. It can process large datasets with minimal cluster resources. XGBoost incorporates a regularized model to prevent overfitting. \cite{chen2016}

\subsection{Multi-Label Classification Algorithms}\label{subsec:multilabel-classification}
\textbf{Multi-label classification} is a supervised learning paradigm where each object is represented by a single instance but is associated with a set of labels, unlike traditional supervised learning, where each example has only one label. The task in multi-label learning is to learn a function to predict the proper label sets for unseen instances. The primary challenge in multi-label learning is the overwhelming size of the output space, as the number of possible label sets grows exponentially with the number of class labels. \cite{zhang2014}

The key to tackling multi-label data lies in exploiting the correlations among labels. Different strategies consider different orders of correlations:
\begin{itemize}
    \item \textit{The first-order strategy} tackles the multi-label learning task in a label-by-label style, ignoring other labels' co-existence. This strategy is beneficial because it is simple and efficient. However, the resulting approach might not be as practical as other strategies due to the ignorance of label correlations.
    \item \textit{Second-order strategy} tackles the task of multi-label learning by considering pairwise relations between labels. Because label correlations are exploited to some extent, the resulting approaches can achieve good generalization performance.
    \item \textit{High-order strategy} tackles the multi-label learning task by considering high-order relations among labels. Although this strategy has more substantial correlation-modeling capabilities than first and second-order strategies, it is more computationally demanding and less scalable.
\end{itemize}

\textbf{Binary Relevance (BR)} is a problem transformation method used in multi-label learning that decomposes a multi-label learning problem into independent binary classification problems. Each classification problem corresponds to a possible label in the label space. BR is a first-order approach that builds classifiers for each label separately and offers the opportunity for parallel implementation. A prominent advantage of BR is its straightforward way of handling multi-label data, which has been employed as the building block of many state-of-the-art multi-label learning techniques. On the other hand, BR completely ignores potential correlations among labels, and the binary classifier for each label may suffer. \cite{zhang2014}

\textbf{Calibrated Label Ranking (CLR)} is a problem transformation method that transforms the multi-label learning problem into a \textit{label ranking problem}, where ranking among labels is achieved by pairwise comparison techniques. It is considered a second-order strategy because it builds classifiers to exploit label correlations for any pair of class labels. CLR incorporates a virtual label into each multi-label training example to specify a thresholding function to bipartition the list of ranked labels into relevant and irrelevant label sets. So, it constructs binary classifiers in a one-vs-one manner, except for those for the virtual label, which mitigates the negative influence of the class imbalance issue. It is worth noting that the training set employed by CLR is identical to the training set employed by Binary Relevance. Therefore, CLR can be regarded as an enriched version of pairwise comparison, where the routine binary classifiers are enlarged with the binary classifiers of BR to facilitate learning. \cite{zhang2014}

\textbf{Classifier Chains (CC)} is a problem transformation method in multi-label learning that transforms the multi-label learning problem into a chain of binary classification problems, where subsequent binary classifiers in the chain are built upon the predictions of preceding ones. It is considered a high-order approach that randomly considers correlations among labels. Compared to BR, CC has the advantage of exploiting label correlations; however, it loses the opportunity of parallel implementation due to its chaining property. \cite{zhang2014}

\textbf{Label Powerset (LP)}, also sometimes referred to as Label Combination (LC), is a problem transformation method used in multi-label classification (MLC). It transforms the multi-label classification problem into a multi-class classification problem. LP operates by treating each unique label set as a separate class. Any classifier suitable for solving a multi-class classification can be applied to solve the newly created single-target multi-class problem. As an advantage, LP preserves label relationships, but it is confined to predicting label sets appearing in the training set, i.e., unable to predict novel label combinations. \cite{zhang2014}

\textbf{Pruned Sets (PSt)} is a problem transformation method that aims to reduce the number of unique classes (label-sets) appearing when a multi-label problem is approached with LP. The process has two phases to achieve this goal. The first, the \textit{pruning step}, removes the infrequently occurring label sets from the training data. The decision on what constitutes an infrequent label-set is a hyperparameter of the method. The method introduces this as a pruning hyperparameter, defining the maximal number of frequent label sets to subsample from the infrequent ones. LP is trained on such a newly created dataset. Another improvement of the method is introducing a threshold function that enables new label combinations to be produced at prediction time. This method's efficiency is advantageous since it is much faster than Label Powerset. On the other hand, the assumptions can break, and the method (without a threshold parameter) cannot introduce novel label sets.

%A \textbf{Conditional Dependency Network (CDN)} is a multi-label classification method aiming to encapsulate the dependencies between labels using dependency networks, which are cyclic-directed graphical models. In these networks, the parents of each variable constitute its Markov blanket. The Markov blanket of a node is the only knowledge needed to predict the behavior of that node and its children. The probabilistic distributions are modeled via simple binary classifier models. These models take the whole feature space augmented by all other labels (excluding the modeled label) as input. While this model preserves label dependencies, the inference performed by Gibbs sampling may require a longer convergence time and may not achieve stationarity if there are many labels. \cite{bogatinovski2022}

\textbf{Meta Binary Relevance (MBR)}, also known as the \textbf{2BR method}, is a method that applies BR in two consecutive stages. In the first stage of the model, $\left| \mathcal{L} \right|$ binary base models are built, where $\mathcal{L}$ refers to the label set. This stage is identical to the standard BR approach, where the multi-label problem is transformed into multiple binary classification problems. In the second stage, the predictions from the first stage augment the feature space; $\left| \mathcal{L} \right|$ features are added. New $\left| \mathcal{L} \right|$ binary models are trained as in BR. \cite{bogatinovski2022}

\textbf{Random $k$ Labelsets (RAkEL)} is an ensemble method for MLC that uses multiple LP models trained on random partitions of the label space. This approach contrasts with training a single LP model on all label combinations. Each of the LP methods learns $2^k$ classes instead of $2^{\left|\mathcal{L}\right|}$, where $k \ll \left|\mathcal{L}\right|$, resulting in multi-class problems that are better balanced. In the evaluation phase, the predictions are obtained by voting among the LP models in the ensemble. Unlike LP, RAkEL uses fewer classifiers than BR, can provide better generalization, and avoids underfitting. However, it does not scale well over time as the number of labels and instances increases.

\textbf{Hierarchy of Multi-label Classifiers (HOMER)} is an ensemble method that transforms the MLC problem into a tree-shaped hierarchy of more straightforward, better-balanced MLC problems, utilizing a divide and conquer strategy. The tree is constructed so the leaves contain singleton labels, while the internal nodes represent joint label sets. In particular, a node includes a training sample if and only if the sample is annotated with at least one of the labels of the label set contained in the node. There are two steps in training a HOMER. First, the tree is built by distributing the labels from the parent node to the children nodes using a balanced clustering algorithm. Then, a multi-label model is trained on a reduced label subset, and the process is repeated until all nodes contain only one label. HOMER is helpful in tasks with many labels where it has demonstrated the best predictive performance. However, the constructed hierarchy is not utilized in problems with a smaller number of labels.


\section{Related Works}\label{sec:related-works}
The objectives of this work can't be achieved without explaining the social media content filtering system. \cite{cotov2023} describes how social media users were selected for content collection. These users included institutions, experts, and hobbyists discussing various topics on social media platforms. Some users focused exclusively on cybersecurity, while others shared a mix of information. When data are divided into training and test sets, selecting an equal number of posts from each user helps understand the author's context. This method tests how well the model can capture the author's content patterns. It relies on the observation that accounts posting about cybersecurity tend to share similar content in the future.

Various methods and models were explored to find the best fit for the dataset to define the filter. The choice of embedding models was crucial because understanding sentence meaning is key to building an effective system in natural language processing.

This filtering system improves real-time cybersecurity awareness by delivering visual alerts through a system-defined interface. The CS-Aware-Next project proposes an architecture that requires components to connect to a common cybersecurity situational awareness system. Data is collected locally and partly in a shared system, creating a split environment. The visualization module handles the user's manual response and communicates with the Self-Healing module.

A similar study by \cite{behzadan2018} introduced a framework for detecting and classifying cyber threat indicators on Twitter using a data-driven approach. This framework employs a Convolutional Neural Network with a binary classifier to identify cyber-related social media posts and a multi-class model to categorize these contents into different cyber threat types.

\cite{behzadan2018} emphasizes the importance of Cyber Threat Intelligence (CTI) for proactive cyber defense, with Open-Source Intelligence (OSINT) from platforms like Twitter being a valuable information source. Many current methods depend on manual rules and keyword-based filtering, which can lead to many false positives and may overlook important information due to flexible language and evolving terms. There is also a lack of open datasets for validation and comparison.

\cite{hussein2023} investigated using machine learning to classify Twitter posts into different cyber threat categories. Their best model used the Random Forest algorithm and achieved an F1 Score of 81{\%}. \cite{masahiroetal2022} focused on machine learning for the multi-label classification of cybersecurity texts to reduce the cost of threat analysis and incident response. Their study reached a maximum micro F1 Score of 71.10{\%}.


\section{Research Gaps}\label{sec:research-gaps}
This thesis addresses several significant research gaps in cybersecurity threat detection and classification from social media data. These gaps are summarized as follows:

\begin{enumerate}
    \item \textbf{Content Relevance Filtering for Cybersecurity Awareness}. Social media platforms serve as vast repositories of cybersecurity information, but the signal-to-noise ratio presents a significant challenge. This research addresses the critical need for intelligent filtering mechanisms to identify relevant cybersecurity content from the overwhelming volume of social media posts. By developing classifiers specifically designed to extract high-value cybersecurity intelligence from general social media content, this work enables organizations to focus on actionable information, enhancing their cybersecurity awareness and response capabilities.
    \item \textbf{Integrated Multi-Label Classification Approaches}. While existing research has primarily focused on binary classification (cybersecurity-related vs.\ non-related) or simple multi-class classification, this thesis explores comprehensive multi-label classification architectures that can simultaneously identify and categorize various cybersecurity threats. This integrated approach is essential for real-world applications where threats often span multiple categories.
    \item \textbf{Advanced Embedding Techniques for Cybersecurity Domain}. Much of the current literature relies on general-purpose embedding models not optimized for cybersecurity terminology and context. This research evaluates specialized embeddings like ATT{\&}CK BERT against general embeddings to determine their effectiveness in capturing domain-specific semantics related to cybersecurity threats.
    \item \textbf{Architectural Comparison}. The thesis presents and compares two architectural approaches, the Two-Step System and the One-Step System, which differ in how the relevance decision is separated from the categorization. The comparison is carried out end to end with paired confidence intervals, so that the trade-offs between modular sequential processing and integrated classification are assessed on the same test posts.
    \item \textbf{Organizational Cybersecurity Awareness}. The research addresses the growing need for organizations to maintain continuous awareness of emerging threats through automated monitoring of social media channels—a vital capability as cyber threats evolve rapidly and information sharing accelerates.
    \item \textbf{Bridging Academic Research and Practical Implementation}. This thesis bridges the gap between theoretical classification models and practical, deployable systems that can enhance organizational cybersecurity posture by aligning with the CS-Aware-Next project's objectives.
    \item \textbf{Adaptability to Evolving Threat Landscapes}. Existing solutions often struggle to adapt to new cybersecurity terminology and emerging threats. This research explores models that can better generalize to novel threats through improved semantic understanding rather than relying solely on keyword matching.
\end{enumerate}

By addressing these research gaps, this thesis contributes to the broader field of applied artificial intelligence in cybersecurity, providing methodologies and an evaluation protocol that can help organizations develop automated threat detection capabilities from social media.
